\scanpage{40}
\deleted{
\begin{flushright}7.4.1967\end{flushright}

Cher Dieudonné,

D'accord qu'il est inutile de rajouter un (vi bis) en Errata à 12.1.1. ; il y aurait lieu de le faire cependant en réédition, car le lecteur trouvera plus \unsure{facilement} dans 12.1.1. (où il regardera tout naturellement) que dans 12.1.3. \unsure{Il|Ia} n'y a pas de raison d'avoir mis l'énoncé sur la prof. rectifiée à part des autres. Certainement tu te canules dans 15.5.6.1. car il est clair à priori que si l'énoncé est \unsure{vrai|vra} lorsque tout fermé irréductible a un \deleted{seul} point irréductible unique, il est vrai aussi sans \unsure{la|fa} restriction d'unicité: il suffit de regarder l'espace \inserted{$X'$} formé des parties fermées irréductibles et l'application $X\quad X'$ qui induit un isomorphisme sur les \deleted{\unsure{tat}} ensembles ordonnés des ouverts de $X$ et $X'$ $\ldots$ Il me semble certain par suite qu'en changeant une ligne ou deux dans ta démonstration, elle doit marcher dans le cas général.
Dans 15.6.3. on peut éviter l'hypothèse noethérienne, car la réduction au cas d'un anneau de valuation discrète est bien inutile, on se ramène à vérifier le relèvement des générisations, qui résulte de 15.6.2.
\unsure{D',accord|D'accord} par contre que la démonstration donnée de 15.6.4. utilise effec-tivement l'hypothèse noethérienne; on l'\unsure{élkiminera|éliminera} dans la réédition, en faisant une réduction au cas noethérien.

As-tu fait le nécessaire pour qu'on m'envoie la suite des 2èmes. épreuves ?

Bien cordialement
}

\mspage{39}
